\documentclass[a4paper,9pt]{extarticle}

%=========================================================
% PACKAGES
%=========================================================

\usepackage[
a4paper,
top=1.8cm,
bottom=1.5cm,
footskip=7mm,
left=1cm,
right=1cm,
headheight=14pt,
headsep=5mm
]{geometry}
\usepackage[T1]{fontenc}
\usepackage[utf8]{inputenc}
\usepackage[french]{babel}

\usepackage[sfdefault]{FiraSans}
\renewcommand{\familydefault}{\sfdefault}

\usepackage{multicol}
\setlength{\columnsep}{0.7cm}

\usepackage{amsmath,amssymb}

\usepackage{titlesec}
\usepackage{enumitem}

\usepackage{tcolorbox}
\tcbuselibrary{skins,breakable}

\usepackage{xcolor}

\usepackage{graphicx}
\usepackage{float}

\usepackage{tikz}

\usepackage{fancyhdr}
\usepackage[hidelinks]{hyperref}


%=========================================================
% PAGE SETTINGS
%=========================================================

\setlength{\parindent}{0pt}
\setlength{\parskip}{1.5pt}

\setlength{\abovedisplayskip}{4pt}
\setlength{\belowdisplayskip}{4pt}

\setlist[itemize]{
leftmargin=*,
itemsep=1pt,
topsep=1pt
}


%=========================================================
% COLORS
%=========================================================

\definecolor{Blue}{HTML}{E3F0FF}
\definecolor{Green}{HTML}{E3F7EA}
\definecolor{Pink}{HTML}{FFE5ED}
\definecolor{Orange}{HTML}{FFF0DC}
\definecolor{Purple}{HTML}{EEE7FF}
\definecolor{Yellow}{HTML}{FFFADB}


\definecolor{BlueTitle}{HTML}{285B8F}
\definecolor{GreenTitle}{HTML}{287052}
\definecolor{PinkTitle}{HTML}{A84063}
\definecolor{OrangeTitle}{HTML}{A96500}
\definecolor{PurpleTitle}{HTML}{6245A5}


%=========================================================
% HEADER
%=========================================================

\pagestyle{fancy}

\fancyhf{}

\fancyhead[L]{\small\textbf{Responsible Software}}
\fancyhead[C]{\small Introduction: Ethics and Stakeholders}
\fancyhead[R]{\small Week 02 --- 14-20/09/2026}

\fancyfoot[C]{\small\href{https://github.com/AlphaXT38G/notes_ba5}{github.com/AlphaXT38G/notes\_ba5}}
\fancyfoot[R]{\small \thepage}

\renewcommand{\headrulewidth}{0.4pt}



%=========================================================
% TITLES
%=========================================================

\titleformat{\section}
{\large\bfseries\color{BlueTitle}}
{}{0pt}{}

\titlespacing{\section}
{0pt}{6pt}{3pt}


% separator

\newcommand{\separator}{
\vspace{2pt}
\hrule
\vspace{3pt}
}


%=========================================================
% BOXES
%=========================================================

% Boxes are kept whole (unbreakable): inside multicols, a box split by tcolorbox
% leaves an empty column behind it.

% DEFINITION - blue
\newtcolorbox{definition}{
enhanced,
unbreakable,
colback=Blue,
colframe=BlueTitle,
boxrule=0.7pt,
arc=0mm,
title=\textbf{Definition},
fonttitle=\bfseries
}

% IMPORTANT - pink
\newtcolorbox{important}{
enhanced,
unbreakable,
colback=Pink,
colframe=PinkTitle,
boxrule=0.7pt,
arc=0mm,
title=\textbf{Key idea},
fonttitle=\bfseries
}

% FORMULAS - green
\newtcolorbox{formula}{
enhanced,
unbreakable,
colback=Green,
colframe=GreenTitle,
boxrule=0.7pt,
arc=0mm,
title=\textbf{Key formulas},
fonttitle=\bfseries
}

% WARNING - orange
\newtcolorbox{warning}{
enhanced,
unbreakable,
colback=Orange,
colframe=OrangeTitle,
boxrule=0.7pt,
arc=0mm,
title=\textbf{Warning},
fonttitle=\bfseries
}

% EXAMPLE - purple
\newtcolorbox{example}{
enhanced,
unbreakable,
colback=Purple,
colframe=PurpleTitle,
boxrule=0.7pt,
arc=0mm,
title=\textbf{Example},
fonttitle=\bfseries
}

% SUMMARY - orange
\newtcolorbox{summary}{
enhanced,
unbreakable,
colback=Yellow,
colframe=OrangeTitle,
boxrule=0.7pt,
arc=0mm,
title=\textbf{Summary},
fonttitle=\bfseries
}



% GOAT - brown: the idea told as simply as possible
\definecolor{Goat}{HTML}{F7EEDF}
\definecolor{GoatTitle}{HTML}{8A5A1F}
\newtcolorbox{goat}{
enhanced,
unbreakable,
colback=Goat,
colframe=GoatTitle,
boxrule=0.7pt,
arc=0mm,
title=\textbf{Explained to a goat},
fonttitle=\bfseries
}



%=========================================================
% FIGURE BOX
%=========================================================

\newtcolorbox{figurebox}{
enhanced,
colback=white,
colframe=gray,
boxrule=0.5pt,
arc=0mm,
center
}



%=========================================================
% DOCUMENT
%=========================================================

\begin{document}

%\footnotesize   --> To reduce the size of the text

\thispagestyle{fancy}



\begin{multicols}{2}


%=========================================================
\section{Responsible Software}


\begin{definition}

\textbf{RESPONSIBLE SOFTWARE} (van de Poel, 2011)

Software \textbf{engineered in a responsible way}: about the \textbf{process}, with a focus on
\textbf{human}, \textbf{active} (forward-looking) responsibility.

\end{definition}


\begin{definition}

\textbf{RESPONSIBLE ENGINEERING}

Making design \textbf{decisions} responsibly: (1) goal to \textbf{do good}; (2)
\textbf{prevent avoidable negative impacts}; (3) take into account people, other systems,
social structures and the planet.

\end{definition}


\begin{definition}

\textbf{ETHICS}

Principles and standards that guide human behaviour.

\begin{itemize}

\item \textbf{Metaethics}: what does ``right'' mean?

\item \textbf{Normative}: how should people act in general?

\item \textbf{Applied}: how to act ethically in a real situation? $\leftarrow$ \textbf{course}

\end{itemize}

\end{definition}


\begin{important}

Responsibility in engineering: norms and standards (ISO, SIA), obligations (certifications),
codes of conduct (ACM), \textbf{+ responsible design methodologies}.

\textbf{Ethics goes beyond the law}: a car meeting all regulations may still need a recall.

\end{important}


\begin{goat}

Responsible software = making your building decisions so that the software helps and does not
hurt, even where the law says nothing.

\end{goat}



%=========================================================
\section{Ethics in Technical Decisions}


\begin{example}

A form with one first name, one last name, no apostrophe, names that never change:
technical choices that \textbf{exclude people}.

\end{example}


\begin{definition}

\textbf{ETHICAL ISSUE}: a situation that may compromise at least one \textbf{ethical value or
principle}.

\textbf{ETHICAL AGENCY}: capacity to decide and act in line with ethical values.

\end{definition}


\begin{formula}

Ethical decision-making (Rest 1986, Schwartz 2016):

\[
\textbf{Recognize} \rightarrow \textbf{Evaluate} \rightarrow \textbf{Act}
\]

\textbf{Ethical sensitivity} = recognize issues (awareness) + identify their impact (who,
what). Grows with knowledge, experience, habits.

\end{formula}


\textbf{Lenses}: desirability (customer), viability (business), feasibility (technology)
\textbf{+ ethicality} (people and environment).


\begin{goat}

Even a name field can be unfair. Step 1 is to notice it; then think; then act.

\end{goat}



%=========================================================
\section{Ethical Dilemmas}


\begin{definition}

\textbf{ETHICAL DILEMMA}

A choice between alternatives that \textbf{all} conflict with ethical values, none overriding
the others: a \textbf{wrong-wrong choice} (trolley problem).

\end{definition}


\begin{warning}

\begin{itemize}

\item \textbf{Technical issue}: no value at stake.

\item \textbf{Ethical issue}: one value compromised.

\item \textbf{Dilemma}: values in conflict whatever you choose.

\end{itemize}

\end{warning}


\begin{example}

\begin{itemize}

\item Brain-to-speech must not capture private thoughts: \textbf{ethical issue} (privacy).

\item ML for medicines also finds chemical weapons: \textbf{dilemma} (\textbf{dual use}).

\item Chatbot conversations sent to teachers to catch hallucinations: \textbf{dilemma}
(learning safety vs. privacy / power).

\end{itemize}

\end{example}


\begin{definition}

\textbf{DUAL USE}: technologies usable for beneficial and harmful purposes. Develop (benefits,
risk of harm) or not (no harm, lost benefits)? GenAI is dual use.

\end{definition}


\textbf{EVIL heuristic} (Lennerfors 2021): \textbf{E}xit, \textbf{V}oice,
\textbf{L}oyalty + \textbf{I}nsubordination.


\begin{definition}

\textbf{JUDGMENT} (Evaluate step): cognitive (reason) + affective (emotion, intuition).

\textbf{MORAL EMOTIONS} = risk indicators. Self: guilt, shame / pride. Other: anger,
contempt, disgust / gratitude, elevation.

\end{definition}


\begin{formula}

\textbf{Normative theories}

\begin{itemize}

\item \textbf{Utilitarianism}: consequences, greatest good for the greatest number.

\item \textbf{Deontology}: the action, universal principles (``never lie'').

\item \textbf{Virtue}: the person (``a courageous, trustworthy person would report it'').

\item \textbf{Care}: relationships, mutual responsibility.

\end{itemize}

\end{formula}


\textbf{Markkula framework}: 1 identify issues, 2 get facts, 3 evaluate options with
\textbf{several} theories, 4 choose and test, 5 implement and reflect.


\begin{goat}

A dilemma: every door hurts someone. Pick the least bad one with good reasons (consequences,
rules, good character, relationships), and listen to your gut alarms.

\end{goat}



%=========================================================
\section{Strategy: Stakeholder Analysis}


\begin{definition}

\textbf{STAKEHOLDER} (Value Sensitive Design)

``Individuals, groups, organizations, institutions, and societies that might reasonably be
affected by the technology'': positive or negative, direct or indirect, possibly
\textbf{non-human}.

\end{definition}


\begin{formula}

\textbf{Question}: who is impacted by my software? \textbf{When}: early.

\begin{enumerate}

\item \textbf{Brainstorm}: daily users, geographically / socially close, activities created or
replaced, maintenance, services needed to run, funding, communities of the domain.

\item \textbf{Direct} (in contact with the software) vs. \textbf{indirect} (impacted without
interacting) to classify and find more.

\end{enumerate}

\end{formula}


\begin{example}

``\textbf{Paper supplier (indirect)}: the software replaces paper records, so the hospital's
orders drop: affected without ever using it.''

\end{example}


\begin{warning}

Graded: explain \textbf{how} it interacts / is impacted, \textbf{specific} to the scenario.
Unclear case: state your hypothesis or argue both ways.

\end{warning}


\begin{goat}

List everyone your app touches: those who use it, and those who never touch it but still feel
it.

\end{goat}



%=========================================================
\section{Strategy: Ethical Speculation}


\begin{formula}

\textbf{Escape the mirror} (Black Mirror writers' room)

\textbf{Part I, dark}: a character hurt by the software, a pitch, a catchy title.

\textbf{Part II, positive}: 1-2 ethical issues, immediate and future consequences, a happy
ending.

\end{formula}


\end{multicols}


\end{document}
